\subsection{Modelo y convenciones}
El robot considerado posee tres articulaciones rotacionales: base, hombro y codo. Se utilizan $q=[\theta_1,\theta_2,\theta_3]^T$, metros, radianes, vectores columna y la base como referencia. Las dimensiones publicadas son $l_1=0.04$ m, $l_2=0.29$ m y $l_3=0.24$ m \parencite{ramzy2024}.

La secuencia de desarrollo sigue las clases: método geométrico, matrices homogéneas, Denavit--Hartenberg (DH), inversa, Jacobianos, singularidades, trayectorias y control. El ejemplo antropomórfico de la semana 7 tiene la misma estructura DH; se sustituyen sus dimensiones por las del paper \parencite{course7}. No se trasladan al VSL las dimensiones, límites ni mallas de los robots utilizados como ejemplos docentes.

Con $q=0$, los eslabones móviles se extienden sobre $+x$ a la altura $l_1$. Se define $c_i=\cos\theta_i$, $s_i=\sin\theta_i$, $c_{23}=\cos(\theta_2+\theta_3)$ y $s_{23}=\sin(\theta_2+\theta_3)$.
